\documentclass[conference,letterpaper]{IEEEtran}
\usepackage[T1]{fontenc}
\usepackage{amsmath,amssymb}
\usepackage{cite}
\usepackage{color}
\usepackage[hidelinks]{hyperref}
\hypersetup{
  pdftitle={Adversarial Cooperation},
  pdfauthor={Santiago Restrepo},
  pdfsubject={Research proposal, September 2026},
  pdfkeywords={adversarial cooperation, applied cryptography, zero knowledge,
    multiparty computation, hardware},
  pdfdisplaydoctitle=true
}
\urlstyle{same}

\newenvironment{gamecases}{%
  \par\vspace{0.65\baselineskip}\begingroup
  \setlength{\parindent}{0pt}%
  \setlength{\parskip}{0pt}%
}{\par\endgroup\vspace{0.8\baselineskip}}
\newcommand{\gamecase}[3][0.65\baselineskip]{%
  \par\nointerlineskip
  \hbox{%
    {\color[gray]{0.65}\vrule width 0.4pt}%
    \hskip7pt
    \vbox{%
      \hsize=\dimexpr\linewidth-7.4pt\relax
      \linewidth=\hsize
      \vspace*{#1}%
      \noindent\textbf{#2}\par\nobreak
      \noindent#3\par}}%
  \penalty0}

\title{Adversarial Cooperation}
\author{\IEEEauthorblockN{Santiago Restrepo}
\IEEEauthorblockA{\href{mailto:contact@waajacu.com}{contact@waajacu.com}\quad
\url{https://waajacu.com/}\\
\textit{Research proposal --- September 2026}}}
\date{}

\begin{document}
\maketitle

\begin{abstract}
Adversarial Cooperation seeks solutions to problems
in game theory using cryptographic protocols without a trusted referee,
and hardware for their execution. The aim is cooperation that preserves
private information and the strategic advantages it supports. We seek
resources to develop, analyze, implement, and demonstrate these protocols
and their hardware.
\end{abstract}

\begin{IEEEkeywords}
Applied cryptography, cooperation, zero knowledge, multiparty computation,
hardware.
\end{IEEEkeywords}

\section{The Cooperation Problem}
The central question is: \emph{how can we construct and demonstrate
cryptographic protocols that let adversarial parties cooperate without a
referee, while retaining their strategic advantage?}
A party should be able to take part in a common game or
verify an agreed result without handing its cards, strategy, or other secrets
to an opponent or a trusted intermediary. We find private and public
institutions are not aware of how exposed they are.

The privacy objective is to reveal nothing about private inputs beyond the
agreed outputs and public information. Each protocol must specify its
adversary model, assumptions, and privacy and correctness requirements.

Candidate tools include cryptographic hashes
and transcript hashing; message authentication codes, authenticated encryption,
and digital signatures; commitments \cite{naor1991}, coin flipping, and secret
sharing; oblivious transfer, zero-knowledge proofs \cite{gmr1989}, and multiparty
computation \cite{gmw1987}; verifiable shuffles, threshold cryptography, private
set operations, and verifiable computation; randomness extraction, key
agreement, and password-authenticated key exchange. Selection depends on the
problem's requirements and assumptions. Prior work on replacing mediators in
strategic games also informs this research \cite{dodis2000}.
No new primitive or general solution is claimed; each proposed
construction requires its own analysis.

\section{Games and Cooperation Problems}
The following cases define initial research targets:
\begin{gamecases}
\gamecase[0pt]{Rock--Paper--Scissors}{Fix each player's choice before either
learns the other's move, then verify the completed round's outcome.}
\gamecase{Tic-Tac-Toe}{Demonstrate a private strategy avoids defeat
under specified game conditions, without disclosing move-selection rules.}
\gamecase{Poker}{Investigate dealing and winner verification without a
trusted dealer, without revealing private cards beyond what the agreed
outcome implies.}
\gamecase{Stag Hunt}{Coordinate mutually beneficial action when each
party's willingness to cooperate depends on the other's participation.}
\gamecase{Joint Optimization}{Optimize an agreed objective using private
inputs or constraints, revealing only the agreed result.}
\end{gamecases}
Each case requires its own specification and analysis. Whether more complex
cooperation can be composed from simpler, elemental games and protocols
remains an open question.

\section{From Specification to Hardware}
The goal is to take formally defined protocols through inspectable executable
representations into specialized hardware that participants can independently
inspect and reproduce. Candidate methods include C reference code, Boolean
circuits, register-transfer-level (RTL) descriptions in Verilog or VHDL, and
FPGA evaluation; choices depend on the protocol.

Translation correctness and correspondence between specification, code,
circuits, and devices are research tasks. Prior work on verified high-level
synthesis provides a foundation \cite{herklotz2021}. Evaluation would address
functional behavior, costs, fabrication and device integrity, and physical
leakage or faults. Functional correspondence alone does not establish
protection against such physical attacks.

\section{Resources and Support}
We seek funding, research collaborators, computing resources, FPGA boards,
and measurement equipment. Support would enable dedicated research,
independent cryptographic review, implementation, and hardware experiments.

\bibliographystyle{IEEEtran}
\bibliography{references}
\end{document}
